% 28_body.tex — 산출물 ㉘ 프라이싱·밸류메트릭 결정서 · 유닛이코노믹스 v1 (빈 템플릿 / 완성 예시 공용 본문) — gen_product_templates.py 가 Product_specs.py 에서 생성
% 드라이버: 28_프라이싱UE_v1_빈템플릿.tex (\wbblanktrue) / 28_프라이싱UE_v1_완성예시.tex (\wbblankfalse — 워크북 §X.5 확정 후)
\documentclass[10pt,a4paper]{article}
\usepackage[a4paper,margin=16mm,top=14mm,bottom=20mm]{geometry}
\def\DocNum{28}\def\DocDay{6}\def\DocIdx{3}
\def\DocTitle{프라이싱·밸류메트릭 결정서 · 유닛이코노믹스 v1}
\def\DocSub{Pricing, Value Metric \& Unit Economics v1}
\def\DocVariant{\B{빈 템플릿}{딥게이지 완성 예시}}
\usepackage{wbstyle}
\def\DocCirc{㉘}   % newunicodechar(㉛~㊵ 폴백) 뒤에 정의해야 활성 문자로 읽힌다
\begin{document}
\docheader
 {\Bg{[회사명 · 제품명]}{딥게이지 · GAUGE}}
 {\Bg{[v1.x / D+n]}{v1.0 D+137 → v1.1 2026-12-11 (D+284)}}
 {\Bg{[작성자 — CEO·CPO]}{강민수 · 윤하경\rd{7mm}}}
 {\Bg{[승인자 — 전원]}{전원\rd{7mm}}}

%% ------------------------------------------------------------------ 1
\secbar{1. 밸류메트릭 4안 비교}
\guide{좌석 / 라인 / 검사기록 건수 / 회피 클레임(성과과금). 다섯 기준으로 H/M/L — 고객 가치 정렬, 예측 가능성(고객), 원가 정렬(우리), 판매 용이성, 확장(NDR) 여지.}
{\footnotesize
\begin{tabularx}{\linewidth}{|L{44mm}|Y|Y|Y|Y|}\hline
\hd{기준} & \hd{좌석} & \hd{라인} & \hd{건수} & \hd{성과과금} \\ \hline
\ifwbblank
\rd{9mm}고객 가치 정렬 &  &  &  &  \\ \hline
\rd{9mm}고객의 예측 가능성 &  &  &  &  \\ \hline
\rd{9mm}원가 정렬 &  &  &  &  \\ \hline
\rd{9mm}판매 용이성 &  &  &  &  \\ \hline
\rd{9mm}확장(NDR) 여지 &  &  &  &  \\ \hline
\rd{9mm}종합 &  &  &  &  \\ \hline
\else
고객 가치 정렬 & L — 검사원 수 ≠ 가치 & M — 라인 = 기록이 나는 곳 & H — 건수 = 사용 & H — 2,700만 회피 \\ \hline
고객의 예측 가능성 & H — 사람 수는 안다 & H — 라인 수는 안다 & L — 월마다 다르다 & L — 예측 불가 \\ \hline
원가 정렬 & L — 계정 수와 무관 & M — 사진 수 편차 4.7배 & H — 건수 ∝ 추론비 & L — 원가와 무관 \\ \hline
판매 용이성 & H — 익숙한 단위 & H — 공장이 세는 단위 & M — '건수가 뭐예요?' & L — 측정·귀속 분쟁 \\ \hline
확장(NDR) 여지 & M — 검사원 수 정체 & L~M — 공장당 3.2 & H — 사용이 늘면 & M \\ \hline
종합 & 검사원이 안 쓰게 만드는 단위 & 선택 — 예측·판매 H, 원가 정렬 M을 알고 포기 & 원가와 맞지만 고객이 이해 못 함 — v2 후보 & 측정 불가 \\ \hline
\fi
\end{tabularx}}

%% ------------------------------------------------------------------ 2
\secbar{2. 선택과 근거}
\guide{무엇을 골랐고(정본: 라인당 월 45만 원, 좌석 무제한) 왜인가. 선택의 대가 — 어떤 고객·어떤 세그먼트에서 이 메트릭이 무너지는가 — 를 한 문단.}
\begin{fieldbox}
\ifwbblank
\lines{6}{9.5mm}
\else
라인당 월 45만 원, 좌석 무제한(커널 행동 2). 근거: ① 검사원이 늘어도 비용이 안 늘어야 검사원이 쓴다(좌석 과금은 '계정 아끼기') ② 라인당 전기·중복 입력 39.3만/월 + 8D·통보·클레임 회피 — 클레임 1건 2,700만 = 라인 3.0 × 45만 × 20개월 ③ 공장이 이미 세는 단위(X-05).

선택의 대가: 원가가 라인이 아니라 사진 건수에 비례하므로 사진 수가 표준(4,100)의 4.7배인 고객(동보프레스 19,400)에서 변동원가율 13.4\% → 51.5\%. 그 고객이 정상인지 예외인지 지금은 모른다 — 항목 4.
\fi
\end{fieldbox}

%% ------------------------------------------------------------------ 3
\secbar{3. 유닛이코노믹스 v1 — 표준 공장 vs 헤비 고객}
\guide{월 사진 건수·VLM 단가·추론비·STT·저장 → 변동원가 → 매출(라인 3.0 × 45만) → 공헌이익률. 정본 §3.2 표.}
{\footnotesize
\begin{tabularx}{\linewidth}{|L{40mm}|Y|Y|}\hline
\hd{항목} & \hd{표준 공장} & \hd{헤비 고객 1곳} \\ \hline
\ifwbblank
\rd{8mm}월 사진 건수 &  &  \\ \hline
\rd{8mm}사진 VLM 단가(원) &  &  \\ \hline
\rd{8mm}사진 추론비 &  &  \\ \hline
\rd{8mm}음성 STT &  &  \\ \hline
\rd{8mm}저장·전송 &  &  \\ \hline
\rd{8mm}변동원가 계 &  &  \\ \hline
\rd{8mm}매출(라인 × 단가) &  &  \\ \hline
\rd{8mm}변동원가율 &  &  \\ \hline
\rd{8mm}공헌이익률 &  &  \\ \hline
\else
월 사진 건수 & 4,100 & 19,400 \\ \hline
사진 VLM 단가(원) & 32 & 32 \\ \hline
사진 추론비 & 13.1 & 62.1 \\ \hline
음성 STT & 3.8 & 3.8 \\ \hline
저장·전송 & 1.2 & 3.6 \\ \hline
변동원가 계 & 18.1 & 69.5 \\ \hline
매출(라인 × 단가) & 135(3.0 × 45만) & 135 \\ \hline
변동원가율 & 13.4\% & 51.5\% \\ \hline
공헌이익률 & 86.6\% & 48.5\% \\ \hline
\fi
\end{tabularx}}

%% ------------------------------------------------------------------ 4
\secbar{4. 원가 붕괴 시나리오와 대응}
\guide{헤비 고객(재촬영 습관)·청구서 780만(예상 210만). 초과분의 귀속을 알려면 무엇이 필요한가(고객별 사용량 계측). 공정 사용 조항·건수 상한·경고 임계값.}
{\footnotesize
\begin{tabularx}{\linewidth}{|L{40mm}|Y|}\hline
\hd{항목} & \hd{내용} \\ \hline
\ifwbblank
\rd{10mm}시나리오(무엇이 원가를 올리는가) &  \\ \hline
\rd{10mm}현재 계측 가능한 것 / 불가능한 것 &  \\ \hline
\rd{10mm}공정 사용(fair use) 조항 &  \\ \hline
\rd{10mm}건수 상한·초과 과금 &  \\ \hline
\rd{10mm}경고 임계값과 조치 &  \\ \hline
\else
시나리오(무엇이 원가를 올리는가) & 12월 청구서 780만(예상 210). 초과 570 중 동보프레스 444(78\%) — 원가표 69.5의 6.39배. 서비스별: VLM 612 · STT 41 · 저장 58 · GPU 69. 가설 4 — H1 재촬영 습관 / H2 앱 자동 재시도 루프 / H3 모델 갱신 재추론 배치 / H4 사진 건수 자체 과소 추정 \\ \hline
현재 계측 가능한 것 / 불가능한 것 & 가능: 테넌트별 API 호출 합계·월 사진 수. 불가능: 사진 단위 재촬영 플래그·호출 단위 재시도·배치별 사진 수·라인별 분해 → H1~H4 중 어느 것인지 알 수 없다 \\ \hline
공정 사용(fair use) 조항 & 신규·갱신 계약부터: 라인당 월 6,000건(표준의 1.5배) 포함, 초과분 건당 40원. 기존 3사에 소급하지 않는다 \\ \hline
건수 상한·초과 과금 & 상한 없음(검사원이 촬영을 주저하면 FN이 는다 — ㉗). 초과 과금만. 동보 12월분은 청구하지 않고 계측 결과 공유 조건으로 1월 협의 \\ \hline
경고 임계값과 조치 & 라인당 월 4,500건 → 팀장 + CS 알림 · 6,000건 → 대표 보고. 첫 조치는 계측(D+300 이후 첫 백로그, 0.5 mm) → 1월 데이터로 판별 → 그때 UX·버그·정책·가격 중 하나 \\ \hline
\fi
\end{tabularx}}

%% ------------------------------------------------------------------ 5
\secbar{5. NDR 상한 추정}
\guide{라인 확장 여지(계약 3.0 vs 실사용) · 좌석 무제한의 뜻 · 어디서 확장 매출이 나오는가. NDR > 100\%가 구조적으로 가능한가.}
{\footnotesize
\begin{tabularx}{\linewidth}{|L{40mm}|Y|L{50mm}|}\hline
\hd{항목} & \hd{값·판단} & \hd{근거} \\ \hline
\ifwbblank
\rd{9mm}고객당 확장 가능 라인 &  &  \\ \hline
\rd{9mm}실사용 대비 여지 &  &  \\ \hline
\rd{9mm}확장 매출의 원천 &  &  \\ \hline
\rd{9mm}NDR 상한(추정) &  &  \\ \hline
\else
고객당 확장 가능 라인 & 공장 평균 3.2 − 계약 3.0 = +0.2(7\%) & 정본 §2.1 \\ \hline
실사용 대비 여지 & 계약 3.0 / 실사용 2.4(G-18) — 축소 위험 20\% & 태산정밀 계약 3 / 실사용 2 \\ \hline
확장 매출의 원천 & ① 라인 증설 +0.2 ② Flow 모듈 분리 과금 — 미정, v2 ③ 포털 규격 추가 유료 옵션 & 좌석 무제한은 확장 원천이 아니다 \\ \hline
NDR 상한(추정) & 약 105\% — GRR을 지키는 것이 먼저 & 실사용 2.4 → 3.0이 NDR보다 급하다 \\ \hline
\fi
\end{tabularx}}

%% ------------------------------------------------------------------ 6
\secbar{6. 파일럿 → 유료 전환 조건}
\guide{구축비 분리 원칙(ARR에 넣지 않는다) · 데이터 이용 · 판정 오류 책임 한계(연 계약액의 n\%) · 해지 조건.}
{\footnotesize
\begin{tabularx}{\linewidth}{|L{40mm}|Y|L{50mm}|}\hline
\hd{조건} & \hd{문구} & \hd{근거·리스크} \\ \hline
\ifwbblank
\rd{10mm}구축비 분리 &  &  \\ \hline
\rd{10mm}데이터 이용·학습 활용 &  &  \\ \hline
\rd{10mm}판정 오류 책임 한계 &  &  \\ \hline
\rd{10mm}해지·환불 &  &  \\ \hline
\rd{10mm}가격 인상 조항 &  &  \\ \hline
\else
구축비 분리 & '구축(마스터 매핑·교육·포털 규격) 비용은 일회성으로 별도 청구하며 구독료에 포함되지 않는다.' 300 · 400 · 500 = 1,200만. 회계: 서비스 매출 계정 — ARR·MRR에 넣지 않는다 & a16z 'bookings ≠ ARR'. 리스크: 보고서에 '계약 ARR'이라는 말이 생긴다 \\ \hline
데이터 이용·학습 활용 & '검사 사진·기록은 비식별화 후 모델 개선에 활용할 수 있다. 도면 이미지·1차사 규격 문서는 제외하며 미업로드 모드를 제공한다.' & B-19 · 김도윤 공문. 리스크: 비식별 기준 정의 \\ \hline
판정 오류 책임 한계 & '판정 초안의 오류로 인한 당사 책임은 연 계약액의 10\%(162만 원)를 상한으로 한다.' — ㉗ 항목 6과 동일 & 순SaaS ACV 1,620만 × 10\%. 리스크: 클레임 1건 2,700만과의 격차 \\ \hline
해지·환불 & 12개월 계약, 30일 전 서면 통지, 잔여 50\% 환불. 구축비 환불 없음 & 12개월 미만 이탈 시 CAC 회수 불가 \\ \hline
가격 인상 조항 & 연 1회, 10\% 이내, 90일 전 통지 & 항목 4 공정 사용 조항의 도입 경로 \\ \hline
\fi
\end{tabularx}}

%% ------------------------------------------------------------------ 7
\secbar{7. 가격표 v1 (공개용 한 장)}
\guide{플랜·단가·포함 사항·제외 사항·최소 계약.}
{\footnotesize
\begin{tabularx}{\linewidth}{|L{30mm}|L{30mm}|Y|Y|L{26mm}|}\hline
\hd{플랜} & \hd{단가} & \hd{포함} & \hd{제외} & \hd{최소 계약} \\ \hline
\ifwbblank
\rd{11mm} &  &  &  &  \\ \hline
\rd{11mm} &  &  &  &  \\ \hline
\rd{11mm} &  &  &  &  \\ \hline
\else
Standard & 라인당 월 45만 원 & Inspect(사진·음성 → 기록, 편집·승인, 매핑, 조도 경고) · Flow(NCR 패키지·H사 포털 내보내기·8D D1~D3·판정 로그) · 좌석 무제한 · 사진 라인당 월 6,000건 · 이메일 지원 & 온프렘 · MES 연동 · 자동판정 · 포털 규격 2종 이상 · 라인 전용 항목 구조 — '검토 중'이라 답하지 않는다 & 라인 2 · 12개월 \\ \hline
구축 패키지 & 300~500만(일회성) & 마스터 매핑(≤ 10영업일) · 검사원 교육 2회 · 포털 규격 1종 설정 & 과거 기록 이관 & 계약당 1회 \\ \hline
(온프렘) & 제공하지 않음 & — & — & ㉖ 항목 5 재검토 조건 \\ \hline
\fi
\end{tabularx}}

\end{document}